% $Id: template.tex 11 2007-04-03 22:25:53Z jpeltier $

\documentclass{styling/vgtc} % final (conference style)

\usepackage{mathptmx}
\usepackage{graphicx}
\usepackage{times}
\usepackage{booktabs}
\usepackage{amssymb}
\usepackage{url}

%% Likert scale helper — call immediately after the question text
\newcommand{\likert}{%
  \par\smallskip
  \begin{center}\small
  \begin{tabular}{@{}ccccccc@{}}
    1 & 2 & 3 & 4 & 5 & 6 & 7 \\[-2pt]
    \multicolumn{2}{l}{\scriptsize Strongly Disagree} &
    & {\scriptsize Neutral} &
    & \multicolumn{2}{r}{\scriptsize Strongly Agree}
  \end{tabular}
  \end{center}\smallskip\noindent}

%% Three blank answer lines
\newcommand{\answerlines}{%
  \par\vspace{4pt}%
  \noindent\rule{\linewidth}{0.3pt}\par\vspace{10pt}%
  \noindent\rule{\linewidth}{0.3pt}\par\vspace{10pt}%
  \noindent\rule{\linewidth}{0.3pt}\par\vspace{6pt}}

\onlineid{0}
\vgtccategory{Research}
\vgtcinsertpkg

\title{DreamGuard: Comparing Mixed Reality Safety Systems}

\author{Andrew Scouten\\ %
        \scriptsize{yzb2@txstate.edu}\\ %
        \scriptsize Texas State University
\and Avery R Vanausdal\\ %
        \scriptsize{arv107@txstate.edu}\\ %
        \scriptsize Texas State University}

%% ─────────────────────────────────────────────────────────────
\abstract{
%% ─────────────────────────────────────────────────────────────
Consumer virtual reality (VR) headsets expose users to physical hazards
by occluding their real environment. Current safety systems, including
Meta's Guardian Grid, rely exclusively on \emph{boundary-aware} triggers:
passthrough is activated when the user approaches a pre-defined play-area
boundary, providing no protection against objects that enter the play
space during a session. We introduce DreamGuard, a mixed reality safety
platform that extends the design space to include \emph{object-aware}
triggers using the device's external cameras and depth sensors. We
implement four passthrough techniques on a consumer Meta Quest headset---
(1)~the Meta Guardian Grid as a boundary-aware baseline, (2)~a
forward-facing Window providing always-on directional awareness,
(3)~a Window augmented with YOLOv9 object detection that reveals
passthrough only at detected hazard locations, and (4)~a Depth Bubble
that exposes the real world wherever the Meta Depth API measures
geometry within a threshold distance. We contribute a five-dimensional
design space that positions these techniques along axes of spatial
awareness model, trigger mechanism, aesthetic integration, information
granularity, and temporal behavior. A within-subjects user study ($N$
participants; data collection ongoing) evaluates the four techniques on
perceived safety awareness, presence, disruption, trust, and trigger
comprehension in a controlled VR dungeon environment. Preliminary
instrumentation confirms that all four conditions are deployable and
that objective trigger-latency logging is operational.
} % end of abstract

\CCScatlist{
  \CCScat{H.5.1}{Information Interfaces and Presentation}{Multimedia Information Systems}%
    {Artificial, augmented, and virtual realities};
  \CCScat{H.5.2}{Information Interfaces and Presentation}{User Interfaces}%
    {Evaluation/methodology};
  \CCScat{H.1.2}{Models and Principles}{User/Machine Systems}%
    {Human factors};
  \CCScat{I.3.7}{Computer Graphics}{Three-Dimensional Graphics and Realism}%
    {Virtual reality};
  \CCScat{I.4.8}{Image Processing and Computer Vision}{Scene Analysis}%
    {Object recognition}%
}

%%%%%%%%%%%%%%%%%%%%%%%%%%%%%%%%%%%%%%%%%%%%%%%%%%%%%%%%%%%%%%%%
%%%%%%%%%%%%%%%%%%%%%% START OF THE PAPER %%%%%%%%%%%%%%%%%%%%%%
%%%%%%%%%%%%%%%%%%%%%%%%%%%%%%%%%%%%%%%%%%%%%%%%%%%%%%%%%%%%%%%%

\begin{document}

\maketitle

%% ─────────────────────────────────────────────────────────────
\section{Introduction}
%% ─────────────────────────────────────────────────────────────

The growth of consumer room-scale VR has made physical collisions a
routine safety concern. Headsets occlude the user's entire field of view,
eliminating awareness of walls, furniture, and other people sharing the
physical space. Meta's response---the Guardian Grid---draws a glowing
boundary grid on top of the virtual scene when the user's head
approaches the edge of the configured play area. While effective for
boundary proximity, the Guardian Grid is fundamentally reactive to
\emph{where the user is} relative to a static boundary, not to
\emph{what objects are present} in the play space. A chair pushed into
the centre of the room, a pet crossing the floor, or a person entering
the space are invisible to the Guardian until the user walks within its
fixed perimeter.

Modern mixed reality (MR) headsets provide two sensor modalities that
make richer awareness possible: (i)~outward-facing colour cameras used
for passthrough video, and (ii)~stereo depth sensors that produce a
per-pixel depth map of the real environment. Neither modality is
exploited by existing consumer safety systems. Prior research has
explored decorative boundary transformations~\cite{qin2024brokenwall},
visual metaphors for guardian proximity~\cite{wu2023guardian}, and
AI-assisted depth-based techniques in constrained laboratory
settings~\cite{sobo2024}, but no published work has examined the
\emph{trigger source}---boundary geometry versus detected objects---as
the primary independent design variable on a shipping consumer platform.

This paper makes three contributions:

\begin{enumerate}

  \item A \textbf{five-dimensional design space} (Section~\ref{sec:designspace})
    that positions existing and novel VR safety transitions along axes of
    spatial awareness model, trigger mechanism, aesthetic integration,
    information granularity, and temporal behavior.

  \item \textbf{DreamGuard}, an open-source MR safety platform (Section~\ref{sec:system})
    implementing four passthrough techniques on Meta Quest, including the
    first depth-driven, object-aware passthrough safety technique on a
    consumer headset.\footnote{\url{https://github.com/TXST-CS7389I-Spring-2026-Group-Project/DreamGuard}}

  \item A \textbf{within-subjects user study protocol} (Section~\ref{sec:study})
    and preliminary instrumentation results comparing the four techniques
    on safety awareness, immersion disruption, trust, and trigger
    comprehension.

\end{enumerate}

The study is in progress. This report presents the system design, study
methodology, and initial pilot data. Full results will be reported upon
completion of data collection.


%% ─────────────────────────────────────────────────────────────
\section{Related Work}
%% ─────────────────────────────────────────────────────────────

\subsection{VR Safety and Guardian Systems}

Meta's Guardian Grid is the de facto consumer safety mechanism for
room-scale VR. It renders a perspective-correct grid overlay in the
user's field of view when the head or controllers approach a
pre-configured boundary. The grid is activated entirely by player
position relative to a static perimeter; no real-time sensing of
in-room objects occurs.

Wu et al.~\cite{wu2023guardian} studied how different visual cues for
boundary proximity affect user awareness and presence, establishing that
the aesthetic \emph{form} of boundary visualisation influences perceived
safety without necessarily improving objective hazard detection. Their
work focuses on the \emph{appearance} of boundary-aware systems and does
not vary the trigger mechanism.

Qin et al.~\cite{qin2024brokenwall} introduced BrokenWall, a system
that replaces the guardian grid with metaphorical ``broken'' virtual
geometry at the boundary, making the safety transition feel diegetic
rather than a superimposed HUD. Again, the trigger is boundary proximity;
the contribution is aesthetic integration.

\subsection{AI- and Depth-Assisted Safety}

SOBO~\cite{sobo2024} is the closest prior work to DreamGuard. It
demonstrated depth- and AI-assisted passthrough triggers on a desktop VR
platform (non-standalone), using a depth camera mounted externally to
detect obstacles. SOBO's key finding---that reactive, object-aware
triggers can be more informative than boundary triggers---motivates the
present work. DreamGuard differs in three ways: (1)~it runs on a
standalone consumer headset using only built-in sensors; (2)~the
object-aware detection uses the headset's forward-facing colour cameras
with an on-device neural network; (3)~the depth-based technique uses
the Meta Depth API's per-pixel environment depth map, which provides
360$^\circ$ coverage without additional hardware.

\subsection{Diminished Virtuality and the Reality--Virtuality Continuum}

Milgram et al.'s Reality--Virtuality (RV) Continuum~\cite{milgram1994} frames
mixed reality as a spectrum between fully real and fully virtual
environments. VR safety transitions are, in this framing,
\emph{involuntary movements along the continuum}: the system forces a
partial reveal of physical reality in response to a hazard signal.

The concept of \emph{Diminished Virtuality} (DV) formalises deliberate
reduction of virtual content to reveal the real world. A recent
taxonomy~\cite{jmir2024dv} distinguishes DV techniques by the spatial
extent of the reveal (full, partial, localised), the trigger modality,
and the temporal profile. DreamGuard's design space builds on this
taxonomy by adding trigger \emph{source} as an explicit axis and
grounding each axis in concrete, implemented examples.

\subsection{Diegetic vs. Non-Diegetic Safety Cues}

Film theory distinguishes diegetic elements (those that exist within the
narrative world) from non-diegetic overlays (those visible to the
audience but not to characters). The Guardian Grid is strictly
non-diegetic: it is a heads-up overlay that has no counterpart in the
virtual world the user inhabits. Diegetic safety cues, where the virtual
environment itself \emph{responds} to proximity or object detection,
represent an open design direction. The techniques in this work vary
along this axis: the Window is non-diegetic, the Depth Bubble is
semi-diegetic (it reveals the real world in depth-shaped regions), and
purely diegetic approaches (virtual objects reacting to real hazards)
are discussed as future work.


%% ─────────────────────────────────────────────────────────────
\section{Design Space}
\label{sec:designspace}
%% ─────────────────────────────────────────────────────────────

We characterise VR safety transitions along five dimensions. Table~\ref{tab:designspace}
positions the four DreamGuard techniques and representative prior work.

\begin{table*}[t]
\centering
\caption{Design space of VR passthrough safety transitions. DG = DreamGuard.}
\label{tab:designspace}
\renewcommand{\arraystretch}{1.25}
\small
\begin{tabular}{@{} p{2.2cm} p{2.0cm} p{2.3cm} p{2.2cm} p{2.3cm} p{2.8cm} @{}}
\toprule
\textbf{Technique} &
\textbf{Spatial Awareness} &
\textbf{Trigger Mechanism} &
\textbf{Aesthetic Integration} &
\textbf{Information Granularity} &
\textbf{Temporal Behavior} \\
\midrule
Guardian Grid (Meta)
  & Boundary-aware
  & Play-area proximity
  & Non-diegetic overlay
  & Boundary only (binary)
  & Persistent while near boundary \\

BrokenWall~\cite{qin2024brokenwall}
  & Boundary-aware
  & Play-area proximity
  & Diegetic metaphor
  & Boundary only
  & Persistent while near boundary \\

SOBO~\cite{sobo2024}
  & Object-aware
  & Depth sensor (external)
  & Partial reveal
  & Per-object depth region
  & Reactive with hysteresis \\
\midrule
DG Window
  & Direction-aware
  & Manual / task-trigger
  & Non-diegetic overlay
  & Forward hemisphere
  & Fixed-duration window \\

DG Window + Detection
  & Object-aware
  & YOLO object detection
  & Non-diegetic, localised
  & Per-detection bounding box
  & Reactive (detection timeout) \\

DG Depth Bubble
  & Object-aware
  & Depth threshold (on-device)
  & Semi-diegetic reveal
  & Per-pixel depth
  & Reactive with GPU threshold \\
\bottomrule
\end{tabular}
\end{table*}

\paragraph{Dimension 1 — Spatial Awareness Model.}
\textit{Boundary-aware} systems know only the static play-area perimeter.
\textit{Object-aware} systems sense in-play objects using active perception.
\textit{Direction-aware} systems provide partial awareness of one
hemisphere without object-level sensing. The Guardian Grid and
BrokenWall are boundary-aware. DreamGuard's Window is direction-aware.
Detection and Depth Bubble are object-aware.

\paragraph{Dimension 2 — Trigger Mechanism.}
Triggers vary from static geometry (boundary proximity), through reactive
sensing (camera detection, depth threshold), to manual activation.
The trigger determines \emph{what events} the system can detect: a cat
entering the play space during a session is invisible to any
boundary-aware trigger but detectable by both the YOLO detector (as
``cat'') and the Depth Bubble (as any object closer than the threshold).

\paragraph{Dimension 3 — Aesthetic Integration.}
The Guardian Grid is purely non-diegetic: the grid exists outside the
virtual world. BrokenWall integrates the safety cue into the virtual
environment's geometry. DreamGuard's Depth Bubble is semi-diegetic: it
reveals real geometry in the shape of the depth field, which aligns
spatially with the virtual scene.

\paragraph{Dimension 4 — Information Granularity.}
Coarse systems (Guardian Grid) communicate only ``you are near the
boundary.'' Fine-grained systems (Depth Bubble) provide a
continuous per-pixel depth signal. Detection-based systems provide
intermediate granularity: object class and bounding box.

\paragraph{Dimension 5 — Temporal Behavior.}
Persistent techniques remain active while the triggering condition
holds (Guardian Grid). Fixed-duration techniques activate for a
pre-set interval (DreamGuard Window in study conditions). Reactive
techniques activate, sustain briefly, and clear after a timeout
(Detection: 3-second timeout; Depth Bubble: continuous while depth
threshold is exceeded).


%% ─────────────────────────────────────────────────────────────
\section{System}
\label{sec:system}
%% ─────────────────────────────────────────────────────────────

\subsection{Platform and Architecture}

DreamGuard runs on Meta Quest (tested on Quest 3) using Unity with the
Meta XR SDK, OpenXR, and Unity's Universal Render Pipeline (URP). The
system is implemented as a Unity package
(\texttt{src/dreamguard/unity/}) consumed by a standalone Unity project
(\texttt{playground/unity/DreamGuard/}). All four passthrough techniques
share a common \texttt{IDreamGuardPassthrough} interface and are managed
by a \texttt{DreamGuardPlayer} component that activates and deactivates
technique GameObjects on demand. The full source is available at
\url{https://github.com/TXST-CS7389I-Spring-2026-Group-Project/DreamGuard}.

\subsection{Guardian Grid (Condition~1)}

Meta's default passthrough behaviour is used unmodified: the system
activates a boundary grid when the user's head or a controller
approaches the configured play-area boundary. No DreamGuard code
intervenes in this condition. The Guardian condition establishes the
boundary-aware baseline against which the other techniques are compared.

\subsection{Windowed Passthrough (Condition~2)}

\texttt{DreamGuardWindowedPassthrough} renders a 0.9\,m $\times$ 0.65\,m
passthrough rectangle positioned 1.5\,m ahead of the user's head using
an \texttt{OVRPassthroughLayer} (UserDefined projection,
\texttt{AddSurfaceGeometry} with a tracked quad). The window is always-on
for the 20-second exposure window triggered by the second orb collection.
Surface geometry registration is deferred to the first \texttt{SetEnabled(true)}
call to prevent a full-screen red/green compositor artefact caused by an
active \texttt{UserDefined/Overlay} layer with unpositioned geometry.

\subsection{Detection-Based Passthrough (Condition~3)}

\texttt{DetectionBasedPassthrough} inherits from a base
\texttt{Detection} class that manages the Meta Passthrough Camera Access
API and a YOLOv9 inference worker running via Unity Sentis (on-device
neural network inference). When the YOLO model detects an object, the
detection callback converts the bounding box from model pixel space to
camera-local ray directions, which are uploaded to the
\texttt{DreamGuard/DetectionPassthrough} shader as per-detection
direction vectors. The shader renders a large world-space sphere
surrounding the camera and punches \texttt{alpha\,=\,0} into the
framebuffer within bounding-box regions, revealing the
\texttt{OVRPassthroughLayer} underlay only at those locations. Up to 16
simultaneous detections are supported; a 3-second timeout clears the
reveal window if no detection fires. The bounding box is expanded by a
configurable factor (default 1.3$\times$) to account for YOLO's tight
fitting to object content.

Direction vectors are stored in render-camera-local space and
transformed to world space per frame via \texttt{unity\_CameraToWorld}
in the shader, making the passthrough hole camera-locked and avoiding
the stereo projection mismatch that arises when \texttt{Camera.main\allowbreak{}.WorldToViewportPoint}
(symmetric mono frustum) is used with a headset's asymmetric per-eye
frustums.

\subsection{Depth Bubble (Condition~4)}

\texttt{PassthroughSphere} uses the Meta Depth API
(\texttt{EnvironmentDepthManager}) to access per-pixel environment depth
in real time. The \texttt{Custom/PassthroughSphere} shader renders a
50-metre world-space sphere centred on the camera. For each fragment,
the shader samples the depth texture; fragments where the real-world
surface is closer than the configured threshold receive
\texttt{alpha\,=\,0} (via \texttt{BlendOp RevSub / Blend One Zero}),
revealing the \texttt{OVRPassthroughLayer} underlay. Fragments where
real geometry is farther than the threshold, or absent, remain opaque
virtual. The sphere is kept centred on the camera in \texttt{Update()}.
\texttt{OVRManager.eyeFovPremultipliedAlphaModeEnabled\,=\,false} is
required so the compositor treats the alpha channel as coverage rather
than premultiplied transparency.

The Depth Bubble provides 360$^\circ$ spatial awareness (bounded only by
depth sensor coverage) and requires no explicit object model---any
real-world geometry within the threshold is revealed, including people,
furniture, and pets.

\subsection{Study Instrumentation}

\texttt{StudyLogger} writes a timestamped CSV (\texttt{study.csv}) to
device persistent storage for each session. Logged event types include
\textsc{session\_start/end}, \textsc{room\_enter/complete},
\textsc{room\_halfway} (triggered on second-orb collection),
\textsc{trigger} (passthrough activation, with technique name and
optional metadata), \textsc{intrusion\_mark} (experimenter-logged),
and \textsc{false\_positive}. Companion files record per-frame player
position, headset pose, and controller pose (delta-filtered at
$\epsilon = 10^{-4}$). Logs are retrieved via \texttt{adb pull} after
each session. Trigger latency is computed as
$\Delta t = t_{\textsc{trigger}} - t_{\textsc{room\_halfway}}$.


%% ─────────────────────────────────────────────────────────────
\section{User Study}
\label{sec:study}
%% ─────────────────────────────────────────────────────────────

\subsection{Study Design}

We conduct a within-subjects study comparing the four passthrough
techniques on a Meta Quest 3 headset. Each participant navigates a
virtual dungeon consisting of four sequential rooms, one per condition,
in a fixed order. The fixed order is an acknowledged limitation; see
Section~\ref{sec:limitations}.

\paragraph{Task.}
Each room contains four collectible orbs. The participant must collect
all four orbs to advance to the next room; an arrow UI at the bottom of
the screen indicates the direction of the next uncollected orb. A counter
displays orbs collected and remaining. Collecting the second orb
activates the room's passthrough technique for 20 seconds before
automatically deactivating.

\paragraph{Conditions.}
\begin{enumerate}
  \item \textbf{Room 1 — Guardian Grid}: Meta's default boundary system.
    Unlike the other three rooms, the Guardian is not activated by orb
    collection; it activates automatically if the participant approaches
    the physical play-area boundary. It may not appear at all if the
    participant remains centred in the space.
  \item \textbf{Room 2 — Window}: A forward-facing passthrough rectangle
    (0.9\,m $\times$ 0.65\,m) is displayed for 20 seconds after the second
    orb is collected.
  \item \textbf{Room 3 — Window + Detection}: For 20 seconds after the
    second orb, YOLO object detection activates. Passthrough is revealed
    only at the bounding boxes of detected objects; no detection means no
    passthrough.
  \item \textbf{Room 4 — Depth Bubble}: For 20 seconds after the second
    orb, the Depth Bubble reveals real-world geometry wherever the depth
    threshold is exceeded.
\end{enumerate}

\paragraph{Apparatus.}
Meta Quest 3 headset. Meta Quest Link (USB) is used during development
and pilot testing; deployed builds are sideloaded for participant sessions.


\subsection{Participants}

Participants are recruited from the Texas State University graduate and
undergraduate student population. Inclusion criteria: age 18 or older,
no photosensitive epilepsy, no prior knowledge of the study hypotheses.
Data collection is ongoing. A pilot session (P00) was conducted by a
member of the research team to validate the logging pipeline and task
flow.

\subsection{Measures}

\paragraph{Subjective (per-room Likert, 7-point scale).}
\begin{itemize}
  \item \textbf{A1} Safety awareness: ``I was aware of physical hazards
    in my real environment while using this system.''
  \item \textbf{A3} Presence: ``I felt present and immersed in the
    virtual environment during this condition.''
  \item \textbf{A4} Disruption: ``The safety system disrupted my
    experience of the virtual world.''
  \item \textbf{A5} Trust: ``I would trust this system to show me the
    real world when I actually needed it.''
  \item \textbf{A8b} Novelty: ``This passthrough technique felt new or
    unfamiliar to me.''
\end{itemize}

\paragraph{Comparative (post-session forced ranking).}
Participants rank the four systems on perceived safety (B1), immersion
disruption (B2), and trust (B3), and indicate their overall preference (B4).

\paragraph{Objective.}
Trigger latency $\Delta t = t_{\textsc{trigger}} - t_{\textsc{room\_halfway}}$,
computed per room from \texttt{study.csv}. Room completion time from
\texttt{rooms.csv}.

\paragraph{Qualitative.}
Open-response items (A9) per room, a post-session verbal interview (C1--C5),
and experimenter observation notes.

\subsection{Procedure}

On arrival, participants complete a demographics and VR experience
questionnaire (age, gender, VR frequency, prior Meta Quest use,
prior VR hazard experience, self-reported interruption sensitivity).
The experimenter reads a standardised briefing script (``you will explore
a dungeon with four rooms; each room uses a different safety system'').
Participants then complete the four-room session without pausing between
rooms. After completing all rooms, participants fill out the per-room
questionnaire from memory, complete the post-session rankings, and
participate in a brief verbal interview. Total session time is
approximately 20--30 minutes.

\subsection{Analysis Plan}

For Likert items with $N \geq 6$ participants, we will apply Friedman's
test across the four conditions followed by Wilcoxon signed-rank pairwise
comparisons (Bonferroni-corrected, $\alpha = 0.05$). For forced rankings,
Kendall's $W$ will assess within-sample agreement. Trigger latency will
be compared across rooms using a repeated-measures design. Qualitative
responses will be analysed using thematic analysis.


%% ─────────────────────────────────────────────────────────────
\section{Preliminary Results}
\label{sec:results}
%% ─────────────────────────────────────────────────────────────

\textbf{Data collection is ongoing.} The following results are from a
single pilot session (P00) conducted by a team member to validate the
study apparatus. No inferential analysis is performed on a single
observation; these data confirm that the logging and analysis pipeline
is operational and characterise the system's objective behaviour.

\subsection{Session Overview}

The pilot participant completed all four rooms in 4.6 minutes, collecting
all 16 orbs and generating 5 trigger events. No intrusion marks and no
false positives were logged.

\subsection{Trigger Latency}

\begin{table}[h]
\centering
\caption{Trigger latency for the pilot session (P00). Latency is
  $t_{\textsc{trigger}} - t_{\textsc{room\_halfway}}$. Room~1 latency
  reflects the Guardian's boundary-proximity trigger; rooms~2--4 use
  the orb-collection trigger.}
\label{tab:latency}
\begin{tabular}{@{}lll@{}}
\toprule
\textbf{Room} & \textbf{Technique} & \textbf{Latency (s)} \\
\midrule
1 & Guardian Grid  & 0.002 \\
2 & Window         & 0.002 \\
3 & Window + Det.  & 0.052 \\
4 & Depth Bubble   & 0.004 \\
\bottomrule
\end{tabular}
\end{table}

Table~\ref{tab:latency} shows sub-100\,ms latency for all techniques.
The slightly higher latency in Room~3 (Window + Detection) is consistent
with the Detection system requiring at least one inference frame to
fire the first \textsc{trigger} event after technique activation. Rooms~1,
2, and~4 trigger within 4\,ms, reflecting near-instantaneous response.


\subsection{Subjective Responses (Pilot)}

The pilot participant reported: safety awareness (A1) of 5 for Window
and 6 for Depth Bubble; presence (A3) of 4 (Window), 4 (Detection), and
6 (Depth Bubble); disruption (A4) of 5 (Window), 5 (Detection), and 4
(Depth Bubble); trust (A5) of 4 (Window), 5 (Detection), and 6 (Depth
Bubble). Guardian Grid items were marked optional and not completed.
The participant ranked the Guardian Grid first (most safe) and the
Window last; ranked the Depth Bubble first for least immersion
disruption; and ranked the Guardian Grid first for trust. Preferred
system for everyday use: Guardian Grid.

The participant noted the Depth Bubble as particularly novel and
surprising: ``The depth bubble because it was new to me and cool to see!''

These single-participant observations are not generalisable. They
suggest the instrumentation is collecting data as intended.


%% ─────────────────────────────────────────────────────────────
\section{Discussion}
\label{sec:discussion}
%% ─────────────────────────────────────────────────────────────

\subsection{Object-Aware vs.\ Boundary-Aware Coverage}

The Guardian Grid cannot detect a chair pushed into the centre of the
room, a pet crossing the floor, or a person entering the space without
approaching the configured boundary. The Detection and Depth Bubble
conditions are, in principle, capable of detecting all of these
scenarios. Whether this broader coverage translates to perceived safety
benefit is an open empirical question that the full study is designed to
answer.

\subsection{The Presence--Safety Trade-off}

Every passthrough technique intrudes on the virtual experience to some
degree. The Window intrudes continuously and visually dominates the
forward field of view, but its spatial extent is predictable. The
Detection window fires reactively and at unexpected screen locations,
which may feel more disruptive even if the intrusion is smaller in area.
The Depth Bubble is spatially contiguous with the real environment's
depth structure, which may feel more natural but could also intrude on
large regions if a table or wall is within the depth threshold.

\subsection{Trigger Comprehension}

The five-dimension design space predicts that object-aware triggers
should produce higher trigger comprehension scores than boundary-aware
ones: when a system shows passthrough at the location of a detected
object, the causal relationship between real-world stimulus and
visual response is more transparent. The Guardian Grid's trigger
mechanism (boundary proximity) is invisible when the boundary is not
reached, which may leave users uncertain about when and why it is absent.
We will test this prediction with items A7 (not yet collected) and the
C3 interview question.

\subsection{Limitations}
\label{sec:limitations}

\paragraph{Fixed condition order.}
The four conditions are experienced in the same sequence by all
participants. Order effects, practice effects, and carryover from prior
conditions are confounded with condition effects. Differences between
rooms cannot be claimed as order-independent without counterbalancing.
We report this as a pilot/sequential within-subjects design and note
counterbalancing as future work.

\paragraph{No true intrusion events.}
The study does not introduce scripted physical hazards (e.g., a
confederate entering the space). The second-orb trigger is a
task-defined activation point, not a response to a real hazard. This
limits ecological validity. The \textsc{intrusion\_mark} log event is
available for experimenter-triggered manual logging of any incidental
real-world intrusions.

\paragraph{Detection reliability.}
YOLOv9 inference on Meta Quest 3 introduces latency (one inference
frame, $\approx$50\,ms in the pilot) and is subject to detection failure.
In the pilot session, the detection-based condition triggered correctly.
Systematic evaluation of false negative rate is deferred to the full
study analysis.

\paragraph{Sample size.}
Data collection is ongoing. Results reported here are from a single
pilot participant and do not support inferential claims.


%% ─────────────────────────────────────────────────────────────
\section{Conclusion and Future Work}
\label{sec:conclusion}
%% ─────────────────────────────────────────────────────────────

We have introduced DreamGuard, a mixed reality safety platform that
extends the design space of VR passthrough transitions from
boundary-aware to object-aware triggering. The platform runs on a
consumer Meta Quest headset, implements four qualitatively distinct
passthrough techniques, and is instrumented for a within-subjects user
study. A five-dimensional design space organises the techniques along
axes of spatial awareness model, trigger mechanism, aesthetic
integration, information granularity, and temporal behavior.

Pilot results confirm sub-100\,ms trigger latency across all four
conditions and validate the logging pipeline. Full participant data
are being collected. Upon completion, we will test whether object-aware
techniques provide measurably higher safety awareness and trigger
comprehension than the boundary-aware Guardian baseline, and whether
those benefits come at an acceptable cost to presence.

\paragraph{Future Work.}
Counterbalanced condition ordering to control order effects.
Adaptive depth-threshold calibration based on play-area size.
Diegetic integration: virtual scene elements that \emph{react} to
detected hazards rather than revealing passthrough. Extension to
multi-person social VR scenarios where other users are a primary hazard.


%% ─────────────────────────────────────────────────────────────
%% Figures
%% ─────────────────────────────────────────────────────────────

\begin{figure*}[p]
  \centering
  \includegraphics[width=\textwidth]{images/final_paper/trajectory_per_session.png}
  \caption{Player trajectory for the pilot session (P00), overlaid on a
    top-down dungeon map. Path colour encodes time (purple = early,
    yellow = late). Symbols mark key logged events (room enter, halfway,
    complete, triggers). TODO: describe observable movement patterns
    and any notable deviations.}
  \label{fig:trajectory}
\end{figure*}

\begin{figure*}[p]
  \centering
  \includegraphics[width=\textwidth]{images/final_paper/latency.png}
  \caption{Trigger latency by condition and per room (pilot session,
    $N=1$). TODO: describe trends across conditions when full data are
    available.}
  \label{fig:latency}
\end{figure*}

\begin{figure*}[p]
  \centering
  \includegraphics[width=\textwidth]{images/final_paper/room_durations.png}
  \caption{Room completion time per condition (pilot session, $N=1$).
    TODO: describe whether technique activation appears to affect
    task completion time.}
  \label{fig:roomdurations}
\end{figure*}

\begin{figure*}[p]
  \centering
  \includegraphics[width=\textwidth]{images/final_paper/timeline.png}
  \caption{Event timeline for the pilot session. Each row represents
    a logged event type; vertical position encodes time. TODO: describe
    the pacing of the session and any notable clustering of events.}
  \label{fig:timeline}
\end{figure*}

\begin{figure*}[p]
  \centering
  \includegraphics[width=\textwidth]{images/final_paper/heatmap_by_condition.png}
  \caption{Position heatmap by condition (pilot session, $N=1$). Warmer
    colours indicate longer dwell time. TODO: describe whether spatial
    coverage differs between conditions once full data are available.}
  \label{fig:heatmapbycondition}
\end{figure*}

\begin{figure*}[p]
  \centering
  \includegraphics[width=\textwidth]{images/final_paper/heatmap_combined.png}
  \caption{Combined position heatmap across all sessions. TODO: describe
    overall traversal patterns and any spatial clustering.}
  \label{fig:heatmapcombined}
\end{figure*}

\begin{figure*}[p]
  \centering
  \includegraphics[width=\textwidth]{images/final_paper/minimap_intrusion_triggers.png}
  \caption{World-space positions of intrusion marks and technique
    triggers (pilot session). Arrows connect intrusion to the subsequent
    trigger event. TODO: describe spatial relationship between where
    hazards occurred and where triggers fired.}
  \label{fig:intrusiontriggers}
\end{figure*}

\begin{figure*}[p]
  \centering
  \includegraphics[width=\textwidth]{images/final_paper/minimap_headset_gaze.png}
  \caption{Headset gaze direction (XZ projection) for the pilot session,
    sampled every 30 frames. Arrow direction indicates where the
    participant was looking. TODO: describe whether gaze aligns with
    expected navigation paths and technique activation regions.}
  \label{fig:gaze}
\end{figure*}

%% ─────────────────────────────────────────────────────────────
\clearpage\onecolumn
\appendix
\section{Study Questionnaire}
\label{app:survey}
%% ─────────────────────────────────────────────────────────────

\begin{center}
  {\large\textbf{DreamGuard: Comparing Mixed Reality Safety Systems}}\\[4pt]
  \textit{CS 7389I\,·\,Texas State University}\\[6pt]
  Participant ID:\;\underline{\hspace{3cm}}\qquad Session Date:\;\underline{\hspace{3cm}}
\end{center}

\bigskip

\subsection*{Before You Begin \textnormal{\small(Experimenter reads aloud)}}

\begin{quote}
``You are going to explore a dungeon with four rooms. Each room uses a different safety system that can show you your real physical environment. Once you have finished all four rooms, I will ask you to fill out a questionnaire reflecting on each room. There are no right or wrong answers---we are interested in your honest experience.''
\end{quote}

\subsection*{Demographics / VR Experience \textnormal{\small(Complete once---before session starts)}}

\begin{enumerate}
  \item Age:\;\underline{\hspace{3cm}}
  \item Gender:\;\underline{\hspace{3cm}}
  \item How often do you use VR headsets?\\[4pt]
        $\square$~None\qquad
        $\square$~A few times total\qquad
        $\square$~Monthly\qquad
        $\square$~Weekly or more frequently
  \item Have you ever used a Meta Quest headset before?\quad $\square$~Yes\quad $\square$~No
  \item Have you ever encountered a real physical hazard (bumped a wall, knocked an object) while using VR?\quad $\square$~Yes\quad $\square$~No
  \item I am generally sensitive to interruptions when I am focused on a task.\likert
\end{enumerate}

\bigskip
\subsection*{Part A --- Per-Room Questionnaire \textnormal{\small(Complete after all four rooms; circle one number per item)}}

\subsubsection*{General Task Questions \textnormal{\small(Complete once before the per-room sections)}}

\noindent\textbf{G1.} The maze task felt cognitively demanding.\likert

\noindent\textbf{G2.} I found the virtual environments in the dungeon engaging.\likert

\bigskip
\subsubsection*{Room 1 --- Meta Guardian Grid}

\noindent\textit{Please answer the following questions related to the first room using the Meta Guardian Grid system.}

\smallskip
\noindent\textit{\textbf{Note:} Unlike the other rooms, the Guardian Grid activates automatically and \textbf{may not have appeared at all}.}

\smallskip
\noindent\textbf{A0.} Did the Guardian Grid appear at any point?\quad $\square$~Yes\quad $\square$~No

\paragraph{Safety Awareness}
\noindent\textbf{A1.} I was aware of physical hazards in my real environment while using this system (optional).\likert

\paragraph{Presence \& Disruption}
\noindent\textbf{A3.} I felt present and immersed in the virtual environment during this condition (optional).\likert

\noindent\textbf{A4.} The safety system disrupted my experience of the virtual world (optional).\likert

\paragraph{Trust}
\noindent\textbf{A5.} I would trust this system to show me the real world when I actually needed it (optional).\likert

\paragraph{Novelty}
\noindent\textbf{A8b.} This passthrough technique felt new or unfamiliar to me (optional).\likert

\paragraph{Open Response}
\noindent\textbf{A9.} Describe what you noticed about this safety system. Did anything surprise you?\answerlines

\subsubsection*{Room 2 --- Window}

\noindent\textit{Please answer the following questions related to the second room using the Window passthrough system.}

\paragraph{Safety Awareness}
\noindent\textbf{A1.} I was aware of physical hazards in my real environment while using this system.\likert

\paragraph{Presence \& Disruption}
\noindent\textbf{A3.} I felt present and immersed in the virtual environment during this condition.\likert

\noindent\textbf{A4.} The safety system disrupted my experience of the virtual world.\likert

\paragraph{Trust}
\noindent\textbf{A5.} I would trust this system to show me the real world when I actually needed it.\likert

\paragraph{Novelty}
\noindent\textbf{A8b.} This passthrough technique felt new or unfamiliar to me.\likert

\paragraph{Open Response}
\noindent\textbf{A9.} Describe what you noticed about this safety system. Did anything surprise you?\answerlines

\subsubsection*{Room 3 --- Window + Detection}

\noindent\textit{Please answer the following questions related to the third room using the Window + Detection passthrough system.}

\paragraph{Safety Awareness}
\noindent\textbf{A1.} I was aware of physical hazards in my real environment while using this system.\likert

\paragraph{Presence \& Disruption}
\noindent\textbf{A3.} I felt present and immersed in the virtual environment during this condition.\likert

\noindent\textbf{A4.} The safety system disrupted my experience of the virtual world.\likert

\paragraph{Trust}
\noindent\textbf{A5.} I would trust this system to show me the real world when I actually needed it.\likert

\paragraph{Novelty}
\noindent\textbf{A8b.} This passthrough technique felt new or unfamiliar to me.\likert

\paragraph{Open Response}
\noindent\textbf{A9.} Describe what you noticed about this safety system. Did anything surprise you?\answerlines

\subsubsection*{Room 4 --- Depth Bubble}

\noindent\textit{Please answer the following questions related to the fourth room using the Depth Bubble passthrough system.}

\paragraph{Safety Awareness}
\noindent\textbf{A1.} I was aware of physical hazards in my real environment while using this system.\likert

\paragraph{Presence \& Disruption}
\noindent\textbf{A3.} I felt present and immersed in the virtual environment during this condition.\likert

\noindent\textbf{A4.} The safety system disrupted my experience of the virtual world.\likert

\paragraph{Trust}
\noindent\textbf{A5.} I would trust this system to show me the real world when I actually needed it.\likert

\paragraph{Novelty}
\noindent\textbf{A8b.} This passthrough technique felt new or unfamiliar to me.\likert

\paragraph{Open Response}
\noindent\textbf{A9.} Describe what you noticed about this safety system. Did anything surprise you?\answerlines

\bigskip
\subsection*{Part B --- Post-Session Comparative Ranking \textnormal{\small(Complete after all four rooms)}}

\subsubsection*{Forced Rankings}

\noindent\textbf{B1.} Rank the four systems by how \textbf{safe} they made you feel (1\,=\,safest):

\smallskip
\begin{center}
\begin{tabular}{@{}lcccc@{}}
\toprule
 & \textbf{1st (most safe)} & \textbf{2nd} & \textbf{3rd} & \textbf{4th (least safe)} \\
\midrule
Meta Guardian Grid & $\bigcirc$ & $\bigcirc$ & $\bigcirc$ & $\bigcirc$ \\
Window             & $\bigcirc$ & $\bigcirc$ & $\bigcirc$ & $\bigcirc$ \\
Window + Detection & $\bigcirc$ & $\bigcirc$ & $\bigcirc$ & $\bigcirc$ \\
Depth Bubble       & $\bigcirc$ & $\bigcirc$ & $\bigcirc$ & $\bigcirc$ \\
\bottomrule
\end{tabular}
\end{center}

\medskip
\noindent\textbf{B2.} Rank the four systems by how much they \textbf{disrupted your immersion} (1\,=\,least disruptive):

\smallskip
\begin{center}
\begin{tabular}{@{}lcccc@{}}
\toprule
 & \textbf{1st (least disruptive)} & \textbf{2nd} & \textbf{3rd} & \textbf{4th (most disruptive)} \\
\midrule
Meta Guardian Grid & $\bigcirc$ & $\bigcirc$ & $\bigcirc$ & $\bigcirc$ \\
Window             & $\bigcirc$ & $\bigcirc$ & $\bigcirc$ & $\bigcirc$ \\
Window + Detection & $\bigcirc$ & $\bigcirc$ & $\bigcirc$ & $\bigcirc$ \\
Depth Bubble       & $\bigcirc$ & $\bigcirc$ & $\bigcirc$ & $\bigcirc$ \\
\bottomrule
\end{tabular}
\end{center}

\medskip
\noindent\textbf{B3.} Rank the four systems by how much you \textbf{trusted} them (1\,=\,most trusted):

\smallskip
\begin{center}
\begin{tabular}{@{}lcccc@{}}
\toprule
 & \textbf{1st (most trusted)} & \textbf{2nd} & \textbf{3rd} & \textbf{4th (least trusted)} \\
\midrule
Meta Guardian Grid & $\bigcirc$ & $\bigcirc$ & $\bigcirc$ & $\bigcirc$ \\
Window             & $\bigcirc$ & $\bigcirc$ & $\bigcirc$ & $\bigcirc$ \\
Window + Detection & $\bigcirc$ & $\bigcirc$ & $\bigcirc$ & $\bigcirc$ \\
Depth Bubble       & $\bigcirc$ & $\bigcirc$ & $\bigcirc$ & $\bigcirc$ \\
\bottomrule
\end{tabular}
\end{center}

\bigskip
\subsubsection*{Overall Preference}

\noindent\textbf{B4.} Which system would you \textbf{most want to use} for everyday VR sessions?\\[6pt]
$\square$~Room 1: Meta Guardian Grid\quad
$\square$~Room 2: Window\quad
$\square$~Room 3: Window + Detection\quad
$\square$~Room 4: Depth Bubble\quad
$\square$~None

\smallskip
\noindent Why?\enspace\rule{0.85\linewidth}{0.3pt}

\subsubsection*{Trigger Moments}

\noindent\textbf{B6.} In which room did the safety system feel most \textbf{surprising} or \textbf{unexpected} when it appeared? What made it feel that way?\answerlines

\subsubsection*{Optional Feedback}

\noindent Any other thoughts or comments?\answerlines

\bigskip
\subsection*{Part C --- Verbal Interview \textnormal{\small(Experimenter asks, takes notes below)}}

\noindent\textit{Read each question aloud. Record key phrases verbatim where possible.}

\medskip
\noindent\textbf{C1.} ``Walk me through how each room's system felt different from the others.''\answerlines

\noindent\textbf{C2.} ``Was there any moment where you felt unsafe during a room? What happened?''\answerlines

\noindent\textbf{C3.} ``Did any system feel like it understood what was happening in your real room? What gave you that impression?''\answerlines

\noindent\textbf{C4.} ``The first room used the Guardian Grid---Meta's default boundary system. How did the other three rooms compare to it?''\answerlines

\noindent\textbf{C5.} ``If you could design the ideal VR safety system, what would it do?''\answerlines

\bigskip
\subsection*{Experimenter Notes}

\noindent Session duration:\;\underline{\hspace{2cm}}~min\qquad
Any technical issues:\;\rule{0.6\linewidth}{0.3pt}

\smallskip
\noindent Notable observations during session:\answerlines

\bigskip
\begin{center}\textit{Thank you for participating.}\end{center}

\end{document}
